% Separate post-CP30-build candidate; no change to the frozen paper.
\subsection{A small point field forces both first-layer diagonal projections}

Let \(K\) be complete discretely valued of mixed characteristic
\((0,p)\), with algebraically closed residue field,
\(p\ge5\), \(e_K=p(p-1)\), and \(\mu_p\subset K\).
Put \(q=p^2\), \(N=q-1\), and \(s=pN\).
Use the common tame good field \(L/K\) of degree \(N\).
The groups below are whole potentially supersingular
height-two Barsotti--Tate groups. All comparisons are actual
isomorphisms of torsion Galois modules, and the matrix layers
are the images of those actual representations.

\begin{lemma}[Neither diagonal projection of the first layer can vanish]
\label{lem:B-RD85-diagonal}
\status{Candidate; complete proof supplied, independent audits pending.}
Suppose two groups with isomorphic actual \(p^3\)-torsion
have \(L\)-rational first-level modules, parameters
\[
 B_i=p^2k_i,\qquad 1\le k_i\le p-2,
       \qquad H_1=\mathbf F_pe_1,\quad H_2=\mathbf F_pe_2.
                                                        \tag{B-RD85.1}
\]
Thus the high lines are exchanged. Let
\(W\subset1+pM_2(\mathbf Z/p^3)\) be the common actual
third-level image over \(L\), and let \(W_1\) be its
image under \(1+pA\longmapsto A\pmod p\).
Both linear maps
\[
        W_1\longrightarrow\mathbf F_p,
        \qquad A\longmapsto A_{11},\quad A\longmapsto A_{22}
                                                        \tag{B-RD85.2}
\]
are nonzero. No bound on the full division-field degree
is assumed.
\end{lemma}
\begin{proof}
The actual opposite-quotient proof of
Lemma~\ref{lem:B-RD83-quotients} and its lattice calculation
shows that the \(12\)- and \(21\)-projections of the same
\(W_1\) are nonzero. This part does not use
\(p\equiv1\pmod4\).

Suppose the \(11\)-projection vanishes. It follows that
\(W_1e_1=\mathbf F_pe_2\). Choose corresponding primitive
second-level points \(P_i\) with \(pP_i=e_1\) in the
common first-level module. This choice means the residual
first-level target is high in model 1 and low in model 2.
The actual action modulo \(p^2\) translates such a point
by \(Ae_1\). Hence its orbit has exactly \(p\) elements.
Because all first-level points are \(L\)-rational, its
point field is Galois, and the actual comparison gives
one common field
\[
       E=L(P_1)=L(P_2),\qquad [E:L]=p,
       \qquad\operatorname{Gal}(E/L)=\mathbf F_pe_2
                                                        \tag{B-RD85.3}
\]
as the translation group. This computation uses the
second-level point orbit, not the order of \(W\).

We check the high inverse polygon explicitly. In the
normalized good coordinate the first-level high target
in model 1 has valuation
\(\alpha_1=s-p^2k_1=p(N-pk_1)\). The coefficient of
\(T^p\) in \([p](T)\) has valuation
\((p-1)p^2k_1\), and the coefficient of \(T^q\) is a
unit. The intermediate universal coefficient bounds
are those in the complete inverse-polygon proof of B69.
The intermediate vertex is strictly above the chord
from \((0,\alpha_1)\) to \((q,0)\), since
\[
          p^2k_1>N-pk_1;
       \qquad p(p+1)k_1>N.                           \tag{B-RD85.4}
\]
The linear coefficient also lies above that chord.
All \(q\) inverse points therefore have valuation
\(\alpha_1/q=(N-pk_1)/p\). The low inverse polygon
in model 2 similarly gives valuation \(k_2/p\).
With the integral valuation on the degree-\(p\) field
\(E\), the two coordinate valuations are consequently
\[
               v_E(T_1(P_1))=N-pk_1,
             \qquad v_E(T_2(P_2))=k_2.               \tag{B-RD85.5}
\]
Both are positive and prime to \(p\).

Every nonidentity automorphism of (B-RD85.3) translates
by a nonzero point of \(\mathbf F_pe_2\), which is low
in model 1 and high in model 2. Formal addition preserves
the coordinate valuation of this displacement. Its two
displacement valuations in \(E\) are thus \(qk_1\) and
\(q(N-pk_2)\), respectively. Applying the exact
prime-valuation equality in
Lemma~\ref{lem:B-RD60-break} to the same automorphism
and the two prime valuations in (B-RD85.5) gives
\[
 \begin{split}
  r(g)&=p(p+1)k_1-N,\\
  r(g)&=qN-(p^3+1)k_2.
 \end{split}                                         \tag{B-RD85.6}
\]
Equality would require
\[
         p(p+1)k_1+(p^3+1)k_2=p^4-1.                \tag{B-RD85.7}
\]
Both coefficients are positive. At the largest permitted
values \(k_1=k_2=p-2\), the right side exceeds the left
side by
\[
                     (p+1)(p^2+1)>0.                \tag{B-RD85.8}
\]
This is a contradiction. Interchanging the models and
the two basis vectors excludes a vanishing
\(22\)-projection in the same way.
\end{proof}

\begin{corollary}[Six wild matrix dimensions in the tame residual sector]
\label{cor:B-RD85-tame}
\status{Candidate; complete proof supplied, independent audits pending.}
Suppose \(p\equiv1\pmod4\). At \(e=p(p-1)\), any
different-special-character actual \(p^3\)-congruent
pair with tame reducible residual inertia has
\[
                \dim W_1\ge3,\qquad\dim W_2\ge3,
                        \qquad |W|\ge p^6.          \tag{B-RD85.9}
\]
In particular the remaining tame reducible sector of
the \(p=5,e=20\) problem has wild image of order at
least \(15625\), and still has the necessary parameter
pair \(\{25,75\}\), type A/B up to a common quadratic
twist.
\end{corollary}
\begin{proof}
B77 excludes square finite determinant, and
Corollary~\ref{cor:B-RD84-first-wild} gives exactly the
exchanged-line parameters (B-RD85.1) at every image size.
The tame spectral projectors in
Proposition~\ref{prop:B-RD83-Lie} put
\(E_{12},E_{21}\) in \(W_1\), and also preserve its
diagonal part. The lemma gives a nonzero diagonal
projection, so that diagonal part is nonzero and
independent of the two off-diagonal directions.
Thus \(\dim W_1\ge3\). Actual \(p\)-th powers give
\(W_1\subset W_2\), hence \(\dim W_2\ge3\).
Since the original residual action over \(L\) is trivial,
\(|W|=p^{\dim W_1+\dim W_2}\), proving the result.
The explicit \(p=5\) pattern is the already proved
corollary to B84; it is not a new existence assertion.
\end{proof}

\begin{remark}[Scope of the diagonal recovery]
\label{rem:B-RD85-scope}
\status{Candidate scope; no unrestricted first-wild recovery claim.}
The point-field lemma holds for every \(p\ge5\) under
its stated exchanged-line hypotheses. Only the
six-dimensional matrix consequence uses
\(p\equiv1\pmod4\) to separate opposite tame weights.
The irreducible and nontrivial residual-wild sectors
retain their earlier bounds; no lower bound \(p^6\)
across all residual types is asserted. No image of order
\(p^6\) or greater is excluded or constructed, and the
unrestricted first-wild least-index interval remains
\(\{2,3\}\).

The argument was suggested by the retained parent note
\path{notes/root_B85_point_field_lead.md}. It uses an
actual degree-\(p\) second-level point field, with the
same automorphism evaluated on two different prime-valued
coordinates. It does not estimate an orbit from only
the ranks of evaluations in two congruence layers.
Nor does a zero trace on the terminal matrix layer by
itself force the entire finite-level determinant to
vanish; such an additional inference is not used.
The full B83 actual quotient and lattice argument and
the relevant B69 inverse-polygon and B60 exact-break
proofs were reread. No new original-source reading is
claimed.
\end{remark}
